\begin{proof}
\emph{The formula for $S_{E}$.} By Theorem \claimref{synthesis.quasi-free-deposit-is-vacuum-plus-shot-noise}(i), for a gauge-invariant quasi-free state (Fact \claimref{hawking.imported-quasi-free-states-car-algebra}) $\mathrm{Var}(N_{f}) = \Tr(fEf(1 - E))$. The cosine moment has $f_{C} = \tfrac12(cH + Hc)$ with $c = \tfrac12(\tau_{k} + \tau_{-k})$, $\tau_{\pm k}$ the shifts of momentum, so $f_{C}$ maps the band space at $p$ to those at $p \pm k$ by $\tfrac12\varepsilon_{\pm k}(p)$; $E$ is diagonal in $p$, so only the terms that go out and back contribute, $\mathrm{Var}\,C = \tfrac14\sum_{p}\sum_{\pm}\tr[\varepsilon_{\pm k}(p)E_{p}\varepsilon_{\pm k}(p)^{\dagger}(1 - E_{p \pm k})]$, and the same for the sine moment; dividing by $V$ gives the display (the two signs of $k$ averaged).

\emph{(i)} At $k = 0$, $\varepsilon = H_{p}$ and the integrand is $\tr[HEH(1 - E)] = \|(1 - E)^{1/2}HE^{1/2}\|_{2}^{2} \ge 0$. It vanishes exactly when $(1 - E)^{1/2}HE^{1/2} = 0$, that is when $H$ maps the range of $E$ into the kernel of $1 - E$, the eigenspace $F_{1}$ of eigenvalue $1$. Write the band space as $F_{1} \oplus F_{m} \oplus F_{0}$ by the eigenvalues $1$, strictly between, and $0$ of $E$. The condition is $H(F_{1} \oplus F_{m}) \subset F_{1}$. Since $H$ is Hermitian and preserves $F_{1}$ it preserves $F_{m} \oplus F_{0}$, so $HF_{m} \subset F_{1} \cap (F_{m} \oplus F_{0}) = 0$: the partly filled modes are zero modes of $H$ and $F_{1}$ is invariant. If $H_{p}$ has no kernel, $F_{m} = 0$ and $E_{p}$ is the projection on an invariant subspace of $H_{p}$; conversely such an $E_{p}$ gives zero. A quasi-free state is pure exactly when $E$ is a projection, and invariant under the quadratic time evolution exactly when $E$ commutes with $H$. For a momentum component the one-body form is multiplication by $p_{1}$ at $k = 0$, which commutes with $E$, so the integrand is $p_{1}^{2}\tr[E_{p}(1 - E_{p})]$, zero exactly when $E_{p}$ is a projection.

\emph{(ii)} In $\Cl(3,0)$ the vector $p$ has $p^{2} = |p|^{2}$, so left multiplication by $p$ on a minimal left ideal has eigenvalues $\pm|p|$ with idempotents $P_{\pm}(p)$; the complex structure of the ideal is multiplication by the central pseudoscalar, and the trace over the ideal of left multiplication by an element $A$ without pseudoscalar part is $2\langle A\rangle_{0}$. A Hermitian operator on the ideal with spectrum in $[0, 1]$ is $\tfrac12(\nu + \mu)$ with eigenvalues $\tfrac12(\nu \pm |\mu|)$, so $|\mu| \le \min(\nu, 2 - \nu)$. Then $\tr[pEp(1 - E)] = \tfrac12\langle p(\nu + \mu)p(2 - \nu - \mu)\rangle_{0}$. The scalar parts: $\langle p\nu p(2 - \nu)\rangle_{0} = \nu(2 - \nu)p^{2}$; the two terms with one factor $\mu$ have no scalar part ($p\mu p$ is a vector); and $\langle p\mu p\mu\rangle_{0} = 2(p\cdot\mu)^{2} - p^{2}\mu^{2}$, since $p\mu p = 2(p\cdot\mu)p - p^{2}\mu$. So the trace is $\tfrac12[\nu(2 - \nu)p^{2} - 2(p\cdot\mu)^{2} + p^{2}\mu^{2}] = \tfrac12p^{2}(\nu(2 - \nu) - \mu^{2}) + (p^{2}\mu^{2} - (p\cdot\mu)^{2})$, and $p^{2}\mu^{2} - (p\cdot\mu)^{2} = -(p\wedge\mu)^{2}$. Likewise $\tr[E(1 - E)] = \tfrac12\langle(\nu + \mu)(2 - \nu - \mu)\rangle_{0} = \tfrac12(\nu(2 - \nu) - \mu^{2})$, which is $n_{+}(1 - n_{+}) + n_{-}(1 - n_{-})$ for $n_{\pm} = \tfrac12(\nu \pm |\mu|)$. The three classes follow: for $\mu = (2n - 1)p/|p|$, $\nu = 1$, the first term is $\tfrac12p^{2}(1 - (2n - 1)^{2}) = 2p^{2}n(1 - n)$ and the second vanishes; for $\nu = 1$ and a unit vector $m$ the first vanishes and $|p\wedge m|^{2} = p^{2}(1 - \cos^{2}\vartheta) = 4p^{2}n(1 - n)$ with $n = \tfrac12(1 - \cos\vartheta)$ the weight of $P_{-}(p)$ in the state, a superposition of the two energy levels at each $p$, that is, of the filled level and a pair; for $E_{p} = P_{\pm}(p)$ both vanish. The check verifies the two identities and the three classes in $\Cl(3,0)$ with rational data.

\emph{(iii)} With $E_{p} = P_{-}(p) + \theta(k_{F} - |p|)P_{+}(p)$, $1 - E_{p + k} = \theta(|p + k| - k_{F})P_{+}(p + k)$. Put $r_{1} = |p|$, $r_{2} = |p + k|$, $\vartheta$ the angle between $p$ and $p + k$, and $\varepsilon = \tfrac12(p + (p + k))$, a vector. Since $P_{+}(q)\,\varepsilon\,P_{\pm}(p) = \tfrac12(r_{2} \pm r_{1})P_{+}(q)P_{\pm}(p)$ for $q = p + k$ and $\tr[P_{+}(q)P_{\pm}(p)] = \tfrac12(1 \pm \cos\vartheta)$, the integrand is $\tfrac14(r_{1} + r_{2})^{2}\tfrac12(1 + \cos\vartheta)$ on the lune $r_{1} \le k_{F} < r_{2}$ (particle to particle) plus $\tfrac14(r_{2} - r_{1})^{2}\tfrac12(1 - \cos\vartheta)$ on $r_{2} > k_{F}$ (sea to particle). The second, over all $p$, is the vacuum's integrand; the state's differs from it by minus its restriction to the blocked region $r_{2} \le k_{F}$, a bounded region. In bipolar coordinates $d^{3}p = (2\pi/k)\,r_{1}r_{2}\,dr_{1}dr_{2}$ on $|r_{1} - r_{2}| \le k \le r_{1} + r_{2}$, and $r_{1}r_{2}\cos\vartheta = \tfrac12(r_{1}^{2} + r_{2}^{2} - k^{2})$, so both integrands are polynomials in $(r_{1}, r_{2})$: for $k \le k_{F}$ the lune is $k_{F} - k \le r_{1} \le k_{F}$, $k_{F} \le r_{2} \le r_{1} + k$, and the blocked region is $0 \le r_{2} \le k_{F}$, $|r_{2} - k| \le r_{1} \le r_{2} + k$. The integrals are $2\pi k_{F}^{5}(\tfrac{\kappa}{2} - \tfrac{\kappa^{5}}{320})$ and $2\pi k_{F}^{5}(\tfrac{\kappa^{4}}{60} - \tfrac{\kappa^{5}}{120})$; dividing by $(2\pi)^{3}$ and subtracting gives the display. (For the blocked part with $r_{2} \ge k$: $\int_{-k}^{k}t^{2}(k^{2} - t^{2})\,dt = \tfrac{4}{15}k^{5}$ in $t = r_{1} - r_{2}$, which gives the leading $\kappa^{4}/60$.) The mean density is $\int_{|p| \le k_{F}}|p|\,d^{3}p/(2\pi)^{3} = k_{F}^{4}/8\pi^{2}$. For $g$ independent fields the variances and the densities add, so $P_{\delta} = S_{E}/(g\bar\rho_{1}^{2})$ and $\Delta^{2}_{\delta} = \tfrac{k^{3}}{2\pi^{2}}\cdot\tfrac{8\pi^{2}k}{g\,k_{F}^{4}}(\dots) = \tfrac{4}{g}\kappa^{4}(\dots)$. The check verifies the overlaps in $\Cl(3,0)$ on rational vectors of rational length and the integrals exactly at rational $k$, with the number structure factors quoted for comparison.

\emph{(iv)} Under (C), $\Delta^{2}_{\zeta} = \gamma^{2}(a_{d}H_{d}/k)^{4}\Delta^{2}_{\delta}$, and $\Delta^{2}_{\delta} \propto k^{3 + s}$ gives $n_{s} - 1 = s - 1$. The two values of $\gamma$: the Poisson relation with $4\pi G\rho = \tfrac32H^{2}$ gives a potential $-\tfrac32(a_{d}H_{d}/k)^{2}\delta$; read as the curvature perturbation of the slice it is $\zeta$ up to sign, $\gamma = \tfrac32$ (the energy constraint on a slice whose expansion is uniform); read as the growing solution's Bardeen potential $\Phi$ (Fact \claimref{synthesis.imported-string-gas-thermal-spectrum-and-conversion}) it gives $\mathcal{R} = \tfrac32\Phi$ in the radiation field (Fact \claimref{synthesis.imported-large-scale-potential-and-sachs-wolfe-relations}), $\gamma = \tfrac94$. For the ball, with physical $k_{F}$ and $\kappa = k/(a_{d}k_{F})$, $(a_{d}H_{d}/k)^{4}\kappa^{4} = (H_{d}/k_{F})^{4}$, and $d\ln(1 - \kappa^{3}/30 + \kappa^{4}/96)/d\ln\kappa = -\kappa^{3}/10 + O(\kappa^{4})$. On observed scales $\kappa$ is of order $10^{-26}$ for a sea that carries the density and is followed by a radiation history (argued: $k_{F}$ is of the order of the temperature of a bath of the same density, for which $k_{*}/(a_{d}T_{d}) < 2 \times 10^{-26}$ by the stored bracket), so the correction $\kappa^{3}/10$ to the index is far below any measurable level. The measured index is $-0.0351 \pm 0.0042$ (Fact \claimref{synthesis.imported-planck-2018-lcdm-parameters}) and $0.0351/0.0042$ is between $8.3$ and $8.4$. With $\rho = gk_{F}^{4}/8\pi^{2} = 3H_{d}^{2}\bar M_{P}^{2}$, $\tfrac4g(H_{d}/k_{F})^{4} = \tfrac4g H_{d}^{2}\cdot\tfrac{\rho}{3\bar M_{P}^{2}}\cdot\tfrac{g}{8\pi^{2}\rho} = H_{d}^{2}/(6\pi^{2}\bar M_{P}^{2})$. Labelling a density by the temperature of a thermal bath of that density, $\rho = \tfrac{\pi^{2}}{30}g_{*}T^{4}$, gives $\gamma^{2}\rho/(18\pi^{2}\bar M_{P}^{4}) = \tfrac{g_{*}}{240}(T/\bar M_{P})^{4}$ at $\gamma = \tfrac32$, in which $\pi$ has cancelled; the check evaluates it at $7.3 \times 10^{8}$ GeV and $10^{12}$ GeV with $g_{*} = 106.75$ and $\bar M_{P} = 2.435 \times 10^{18}$ GeV, and brackets the density at which it equals $A_{s} = 2.1 \times 10^{-9}$, $\rho = 8\pi^{2}A_{s}\bar M_{P}^{4}$, with $333/106 < \pi < 355/113$.

\emph{(v)} For $T \ll k_{F}$ and $k \ll k_{F}$ the occupations vary only near the Fermi surface, where the energies in $\varepsilon$ are $k_{F}$ and the overlap is $1$, so $S_{E} = k_{F}^{2}\int n_{p}(1 - n_{p + k})\,d^{3}p/(2\pi)^{3}$ with $n$ the Fermi-Dirac function of $|p| - k_{F}$. With $\mu$ the cosine of the angle between $p$ and $k$, $|p + k| - |p| = k\mu$, and $\int n(\epsilon)(1 - n(\epsilon + \Delta))\,d\epsilon = \Delta/(1 - e^{-\Delta/T})$, since $n(\epsilon)(1 - n(\epsilon + \Delta)) = [n(\epsilon) - n(\epsilon + \Delta)]/(1 - e^{-\Delta/T})$ (an identity of rational functions of $e^{\epsilon/T}$, in the check) and $\int[n(\epsilon) - n(\epsilon + \Delta)]\,d\epsilon = \Delta$. So $S_{E} = \tfrac{k_{F}^{4}}{4\pi^{2}}\int_{-1}^{1}\tfrac{k\mu}{1 - e^{-k\mu/T}}\,d\mu = \tfrac{k_{F}^{4}k}{4\pi^{2}}\int_{0}^{1}\mu\operatorname{coth}\tfrac{k\mu}{2T}\,d\mu$, which is $\tfrac{k_{F}^{4}k}{8\pi^{2}}F(k/T)$. For $y \ll 1$, $\operatorname{coth} z = 1/z + z/3 + \dots$ gives $F = 4/y + y/9$; so $S_{E}(0) = \tfrac{k_{F}^{4}}{2\pi^{2}}T = 4\bar\rho T$, in agreement with (i) ($\int\omega^{2}n(1 - n) = 4T\bar\rho$ for massless quanta). For $y \gg 1$, $F = 1 + 2\int_{0}^{1}\mu(\operatorname{coth}\tfrac{y\mu}{2} - 1)\,d\mu = 1 + \tfrac{8}{y^{2}}\int_{0}^{y/2}z(\operatorname{coth} z - 1)\,dz$, and $\int_{0}^{\infty}z(\operatorname{coth} z - 1)\,dz = 2\sum_{n \ge 1}\int_{0}^{\infty}ze^{-2nz}dz = \tfrac12\zeta(2) = \pi^{2}/12$, the remainder being exponentially small: $F = 1 + 2\pi^{2}/3y^{2}$. Under (C) the index is $d\ln F/d\ln y$: $-1 + y^{2}/18$ at small $y$, and with $e = 2\pi^{2}/3y^{2}$, $de/d\ln y = -2e$, it is $-2e/(1 + e)$ at large $y$, with derivative $4e/(1 + e)^{2}$. Setting $-2e/(1 + e) = -t$ gives $e = t/(2 - t)$, $1 + e = 2/(2 - t)$ and running $t(2 - t)$; at $t = 0.0351$, $e = 0.01786$, $y = (2\pi^{2}/3e)^{1/2}$ between $19.1$ and $19.3$, where the exponential remainder is below $10^{-7}$. For $S = S(0) + ck$ the index is $-f$ with $f = S(0)/(S(0) + ck)$ and $df/d\ln k = -f(1 - f)$, running $f(1 - f)$.

\emph{(vi) (Argued.)} A quasi-free state's entropy is $-\Tr[E\ln E + (1 - E)\ln(1 - E)]$, a function of the spectrum of $E$, which quadratic flows and the number-diagonal generators that (L$_{E}$) permits conserve (Theorem \claimref{synthesis.quadratic-flows-cannot-thermalise-descent}(i), (ii)). So a deposit that becomes today's radiation under those dynamics has the record's entropy at the descent, and the stored relation $a_{d}T_{d}\,g_{s}^{1/3} = a_{0}T_{0}\,g_{s0}^{1/3}$ (Proposition \claimref{synthesis.finite-range-descent-deposit-cannot-seed-cmb}(iii)) with $k_{*}/(a_{0}T_{0}) < 1.37 \times 10^{-27}$ and $g_{s} \le 10^{4}$ gives $k_{*}/(a_{d}T_{d}) < 14 \times 1.37 \times 10^{-27} < 2 \times 10^{-26}$ (the bracket of the stored proof of clause (iv) of Theorem \claimref{synthesis.quasi-free-deposit-is-vacuum-plus-shot-noise}). By (v) the white regime is $k/a_{d} \ll 4T_{d}$; the observed band is inside it by a factor above $10^{26}$. A purification of the occupations by partner modes on the slice is a pure quasi-free state of modes and partners together, and by (i) its white term vanishes exactly when its one-particle density is a projection commuting with the joint one-particle Hamiltonian. A homogeneous state mixes a mode only with partners at the same momentum, and a projection that commutes with the Hamiltonian mixes only levels of equal energy; so the partners are degenerate with their modes, the joint density at $p$ is a projection $E_{p}$ of rank $r_{p} \ge 1$ on the degenerate band space wherever the mode's occupation $n_{p} = \langle 1|E_{p}|1\rangle$ is strictly between $0$ and $1$, and the partners' occupations sum to $r_{p} - n_{p} \ge 1 - n_{p}$. Their energy density is then at least $\int\omega_{p}(1 - n_{p})\,d^{3}p/(2\pi)^{3}$ over the momenta the modes occupy, which diverges for Fermi-Dirac occupations (positive at every momentum). The pair states of Theorem \claimref{synthesis.quadratic-flows-cannot-thermalise-descent}(iii), $c\ket{00} + s\ket{11}$, superpose levels whose energies differ by the sum of the two one-particle energies; with both positive on the slice the state is not stationary, and in the Weyl form of (ii) its white term is $4p^{2}n(1 - n)$ per momentum, twice the mixed state's $2p^{2}n(1 - n)$.
\end{proof}
